% ==============================================================================
% Psychology Paper Template  (strict APA 7th edition structure)
%
% APA Publication Manual (7th ed.); APA 7 statistical reporting (2019):
% effect sizes + 95% CIs + exact p-values; pre-registration via OSF /
% AsPredicted; open-science data / materials / code sharing.
%
% Citation style: APA 7 author-year, approximated by natbib + apalike.bst.
% For strict APA 7, replace with biblatex [style=apa] + biber.
%
% Compile:  pdflatex main -> bibtex main -> pdflatex main -> pdflatex main
% ==============================================================================
\documentclass[11pt,a4paper]{article}

% ---- Page geometry ----
\usepackage[a4paper,top=25mm,bottom=25mm,left=25mm,right=25mm]{geometry}

% ---- Fonts & encoding (APA 7: Times New Roman 11pt or Georgia 11pt) ----
\usepackage[T1]{fontenc}
\usepackage[utf8]{inputenc}
\usepackage{mathptmx}
\usepackage{helvet}
\usepackage{courier}

% ---- Line numbers + double spacing for review (APA 7 manuscript) ----
\usepackage{lineno}
\usepackage{setspace}
\doublespacing                    % APA 7 manuscript: double-spaced throughout
\linenumbers                      % APA 7 submission requires line numbers

% ---- Math, figures, tables ----
\usepackage{amsmath,amssymb}
\usepackage{graphicx}
\usepackage{booktabs}
\usepackage{array}
\usepackage{xcolor}
\usepackage{multirow}

% ---- APA-style section heading formatting (APA 7 Level 1: bold, centered) ----
% The article class defaults are close enough; for strict APA 7 typography
% use the `apa7` document class from CTAN instead.

% ---- Citations: APA-style author-year (approximated) ----
\usepackage[round,authoryear]{natbib}
\bibliographystyle{apalike}       % swap with plainnat or biblatex-apa for strict APA 7

% ---- Hyperref last ----
% \usepackage{hyperref}
% \hypersetup{colorlinks=true,linkcolor=blue,citecolor=blue,urlcolor=blue}

\title{A Concise and Informative Title of the Psychology Submission}
\author{First Author\textsuperscript{1,*}, Second Author\textsuperscript{2}, Senior Author\textsuperscript{1}\\
\textsuperscript{1}Department of Psychology, University of Y, City, Country\\
\textsuperscript{2}Department of Psychological Sciences, Institute of W, City, Country\\
\textsuperscript{*}e-mail: first.author@univ.edu}
\date{}

\begin{document}
\maketitle

% ==============================================================================
% Abstract  (150-250 words; APA 7)
% ==============================================================================
\begin{abstract}
% TODO: 150-250 words. Sentence 1-2: purpose and research question. Sentence
% 3-4: participants (n, demographics) and design. Sentence 5-6: key findings
% with effect sizes (Cohen's $d$ or $\eta^2$) and 95\% CIs. Sentence 7:
% theoretical or practical implications.
\end{abstract}

\noindent\textbf{Keywords:} % TODO: 3-5 keywords separated by commas.
keyword one, keyword two, keyword three

% ==============================================================================
% 1. Introduction  (no "Introduction" heading per APA 7; CARS model)
% ==============================================================================
\section{Introduction}
\label{sec:intro}
% TODO: review prior literature; identify the gap; state the research
% question and the testable hypotheses H1, H2, ... (\texttt{\% TODO:
% hypothesis}). APA 7 does not print an "Introduction" heading.

% ==============================================================================
% 2. Method  (APA 7: subsections Participants / Design / Materials /
% Procedure / Data Analysis)
% ==============================================================================
\section{Method}
\label{sec:method}

\subsection{Participants}
% TODO participants: report $N$, recruitment channel (e.g. university subject
% pool, Prolific, MTurk), inclusion / exclusion criteria, demographic
% breakdown (gender, age $M \pm SD$, ethnicity), and an a-priori power
% analysis justifying the sample size. State IRB approval and informed
% consent procedure here (\texttt{\% TODO: irb}). Disclose any data
% exclusion (e.g. attention-check failures) with the count.

\subsection{Research Design}
% TODO design: state the design (e.g. 2 (Condition: experimental vs.
% control) $\times$ 2 (Time: pre vs. post) mixed-factorial design, with
% Condition as the between-subjects factor and Time as the
% within-subjects factor). List the independent variables and their levels,
% the dependent variable(s) with their scale of measurement, and any
% control variables.

\subsection{Materials}
% TODO materials: list each scale / stimulus with its source citation,
% number of items, scoring, and reliability evidence (Cronbach's $\alpha$
% from the original study AND from the present sample). For experimental
% stimuli, describe selection rationale and counterbalancing. For equipment
% (e.g. eye-tracker, EEG), state the model and manufacturer.

\subsection{Procedure}
% TODO procedure: step-by-step description -- informed consent ->
% instructions -> practice trials -> experimental task -> debriefing.
% State the randomisation method (e.g. complete randomisation via
% Qualtrics), the number of trials, the trial timing, and the
% compensation.

\subsection{Data Analysis}
% TODO analysis: specify the primary analysis model (e.g. mixed ANOVA,
% multilevel model, Bayesian regression), the effect-size metric
% (Cohen's $d$, $\eta_p^2$, $\omega^2$, OR), the 95\% CI method, the
% handling of missing data, the alpha level ($\alpha = .05$), and the
% software + version (e.g. R 4.3.1 with lme4 1.1-34).

% ==============================================================================
% 3. Results
% ==============================================================================
\section{Results}
\label{sec:results}
% TODO: report each hypothesis test in order. APA 7 statistical reporting:
%   - report exact $p$ values to three decimals ($p = .034$), use $p < .001$
%     for very small values; do NOT use $p < .05$.
%   - report effect sizes with 95\% CIs: $d = 0.45$, 95\% CI [0.21, 0.69].
%   - report descriptive statistics: $M = 5.32$, $SD = 1.24$, $n = 120$.
%   - reference each table and figure inline.

\subsection{Preliminary analyses}
% TODO: data screening, assumption checks (normality, homogeneity of
% variance, sphericity), and any data exclusions.

\subsection{Primary analysis}
% TODO: test of the primary hypothesis; reference Table 1 and Figure 1.

\begin{table}[t]
\caption{Descriptive statistics and primary hypothesis tests. Caption ABOVE the table (APA convention).}
\label{tab:primary}
\centering
\begin{tabular}{lcccc}
\toprule
Condition & $n$ & $M$ & $SD$ & Statistic \\
\midrule
Experimental & --- & --- & --- & --- \\
Control      & --- & --- & --- & --- \\
\multicolumn{5}{l}{Cohen's $d$ = ---, 95\% CI [---, ---], $p$ = ---} \\
\bottomrule
\end{tabular}
\end{table}

\subsection{Secondary analyses}
% TODO: exploratory or pre-specified secondary analyses.

\subsection{Sensitivity analyses}
% TODO: re-run the primary model under alternative specifications (e.g.
% intention-to-treat vs per-protocol; alternative missing-data handling;
% Bayesian re-analysis with informed prior).

% ==============================================================================
% 4. Discussion
% ==============================================================================
\section{Discussion}
\label{sec:discussion}
% TODO: (i) restate the key findings without statistics, (ii) interpret in
% light of prior literature and theory, (iii) theoretical contribution,
% (iv) practical implications, (v) limitations
% (\texttt{\% TODO: limitations}), (vi) future research.

\subsection{Limitations and Future Research}
% TODO: sample restrictions, self-report bias, single-method, external
% validity, demand characteristics, etc.

% ==============================================================================
% 5. Conclusion
% ==============================================================================
\section{Conclusion}
\label{sec:conclusion}
% TODO: one paragraph stating the contribution and a forward-looking statement.

% ==============================================================================
% Declarations  (APA 7 + Open Science)
% ==============================================================================
\section*{Declarations}

\noindent\textbf{Funding:} % TODO: grant numbers and funding bodies.

\noindent\textbf{Conflicts of interest:} % TODO: declare or state none.

\noindent\textbf{Ethics approval:} % TODO irb: IRB name + approval number.

\noindent\textbf{Informed consent:} % TODO: state written / verbal consent
% procedure; for minors, state parental consent + child assent.

\noindent\textbf{Pre-registration:} % TODO preregistration: state the
% registry (OSF / AsPredicted) and provide the URL and date.

\noindent\textbf{Open science:} % TODO open-science: provide the OSF / GitHub
% / Zenodo link to data, materials, and analysis code. State any
% restrictions (e.g. clinical data cannot be shared publicly).

% ==============================================================================
% Acknowledgements
% ==============================================================================
\section*{Acknowledgements}
% TODO: funders, research assistants, participants.

% ==============================================================================
% References  —  APA 7 author-year style
% Compile:  pdflatex main -> bibtex main -> pdflatex main -> pdflatex main
% ==============================================================================
\bibliography{references}
% Inline fallback (uncomment if not using BibTeX):
% \begin{thebibliography}{99}
% \bibitem{ref_example} Author, F., \& Author, S. (20YY). Title of the paper.
%   \emph{Journal of Experimental Psychology: General, XX}(Y), 1--10.
%   https://doi.org/XX.XXXX/xxxxxxxx
% \end{thebibliography}

% ==============================================================================
% Appendices  —  full scales, stimuli, supplementary analyses
% ==============================================================================
\appendix

\section{Full Scale Items}
\label{app:scales}
% TODO: list the full wording of each scale item with the response anchors.

\section{Stimuli}
\label{app:stimuli}
% TODO: list the full set of experimental stimuli with counterbalancing.

\section{Supplementary Analyses}
\label{app:supp}
% TODO: Bayesian re-analyses, equivalence tests, per-participant data plots.

\end{document}
